\section{Introduction}\label{sec:intro}

After earthquakes, floods and large fires, cellular and backhaul infrastructure is often unavailable
precisely where responders need information. Small multirotor UAVs carrying thermal cameras can
search such areas quickly, and onboard detectors now run in a few milliseconds on embedded GPUs, so
victims can be detected without streaming video to a server \cite{hituav}. What remains hard is
getting each detection to the responders \emph{in time}. Without infrastructure, a report travels
only over the UAVs' own short-range radio links, which are obstructed by buildings, fade, and break
when UAVs fail.

Two families of methods address the communication side. Connectivity-maintenance controllers keep
the swarm's communication graph connected at all times, either through the algebraic connectivity
of the graph \cite{sabattini2013,zavlanos2008} or through edge potentials that never let a link
break \cite{ji2007}. Every detection then reaches the responder immediately, but the whole swarm is
tied to the responder's neighbourhood, which limits how far it can search. Intermittent-connectivity
and data-ferrying approaches let robots disconnect and meet according to a schedule
\cite{kantaros2017}, or plan ferry routes to sources whose positions are known in advance
\cite{cui2023ferry,bafarassat2026aoi}. In search and rescue neither assumption fits well: the information
sources, the victims, are \emph{discovered by the search itself}, at unknown places and times.

In an earlier version of this system we coupled a thermal detector with particle swarm
optimisation (PSO) and a distance penalty that kept UAVs within 50\,m of their nearest neighbour.
That design has three limitations, which also shaped this work. A fixed distance does not describe
a radio link behind rubble; keeping the nearest neighbour close does not keep the network connected
(Lemma~\ref{lem:nn}); and the evaluation used a 25\,m area that is far smaller than the radio range,
in which, as we show in Section~\ref{sec:eval:sitl}, every method stays connected and the methods
cannot be told apart.

This paper studies the \emph{discovery--delivery trade-off} directly. We measure the share of
victims whose report reaches the responder within a deadline, and we treat connectivity as a means
to that end rather than as an objective. Our contributions are:
\begin{itemize}
\item \textbf{Information-triggered delivery.} Each UAV searches without connectivity constraints and
switches to a delivery mode when a victim it confirmed has not reached the responder after a hold
time $T_h$ (Section~\ref{sec:method}). Rendezvous are triggered by information that the search
itself discovers, rather than by a schedule or by known source locations. $T_h$ exposes the
trade-off explicitly, and we bound the probability that a confirmed detection is lost to UAV
failures before it is secured (Proposition~\ref{prop:exposure}).
\item \textbf{Corroborative multi-view sensing.} Weak detections request a lateral second view from a
neighbour, and views from different directions are fused so that occluded viewpoints cannot veto a
clear one (Section~\ref{sec:fusion}). In dense swarms this fusion rule is the largest single
contributor to performance.
\item \textbf{A link-preservation safety layer with guarantees.} While a delivery is in progress,
the UAVs on the delivery path constrain their steps using a realistic outage model (path loss,
obstructions, shadowing, fading, interference) and a local minimum spanning tree. Preserved links
remain reliable for every feasible step (Proposition~\ref{prop:link}), and the delivering component
stays connected (Corollary~\ref{cor:conn}). We show why nearest-neighbour constraints do not
suffice (Lemma~\ref{lem:nn}).
\item \textbf{A reproducible evaluation.} We compare against seven baselines covering
communication-aware optimisation, algebraic-connectivity control, potential-based topology control,
multi-agent reinforcement learning and classical sweep search, all tuned with the same budget. The
evaluation spans 3--50 UAVs, four channel profiles, six failure models, moving obstacles and
detector uncertainty, uses 30 seeds per configuration with paired significance tests, and is
cross-checked against PX4 software-in-the-loop flight in Gazebo. Code, scenarios and results are
released.
\end{itemize}

Across the scenarios SENTINEL reports 23--30\,\% more victims within 60\,s than the best baseline in
urban rubble under all failure models, and it is significantly better than all baselines on the
400\,m area at every swarm size from 3 to 50; in the harshest channel it is ahead of or tied with the
best baseline. The ablation shows where this comes from: the delivery trigger and the fusion rule
account for the gains, whereas the link-preservation layer has no significant effect on performance
and fulfils its safety role, with $3.4\times10^{-5}$ link violations per preserved link-step. We
also report what does not work as well: keeping the swarm connected at all times lowers on-time
reporting to 0.03--0.13; in small, dense swarms a MAPPO policy trained to convergence comes close to
SENTINEL (differences not significant for $N=10$--20 on the 200\,m area); and the thermal detector,
trained on real imagery, transfers poorly to rendered Gazebo scenes.

Section~\ref{sec:related} positions this work against prior approaches,
Section~\ref{sec:model} defines the system model and problem, Section~\ref{sec:method} presents
SENTINEL, Section~\ref{sec:impl} describes the implementation, Section~\ref{sec:eval} reports the
evaluation and Section~\ref{sec:discussion} discusses limitations.
